% ============================================================
%  Literature Review
%  Covers all 14 verified papers in docs/related_work.md.
%  Numbers that come from the paper's own results are macros;
%  everything else is narrative.
% ============================================================

% ------------------------------------------------------------------
\subsection{Insider-Threat Detection}
% ------------------------------------------------------------------

\textbf{Le and Zincir-Heywood}~\cite{le2021ensembles} combine four unsupervised anomaly
detectors---Isolation Forest, LOF, one-class SVM, and a neural autoencoder---into an
ensemble scored against temporal user-behaviour features derived from the CERT dataset.
The work is notable for framing evaluation as an \emph{investigation budget}: the ensemble
achieves coverage of 60\,\% of malicious insiders when the analyst reviews only 0.1\,\% of all
user-sessions, rising to full coverage below the 5\,\% mark.
Three structural differences separate it from our setting. First, the budget is expressed as a
fraction of the total session pool across the entire evaluation window, rather than as a fixed
daily queue size; translating that fraction into a per-day alert count is not straightforward.
Second, each alert carries no explanation, so an analyst cannot reason about why a session was
flagged. Third, the underlying logs are on-premise host events, whereas our target is the
who-what-where-when structure of cloud audit records.

\textbf{Yuan and Wu}~\cite{yuan2021deep} survey deep-learning approaches to insider-threat
detection across eleven corpora, the majority of which are CERT variants, and observe that the
research community has grown heavily dependent on simulator-generated data.
Three recurring deficiencies are identified: the scarcity of real-world labelled corpora, severe
class imbalance between normal and malicious activity, and a tendency to evaluate using aggregate
metrics such as AUC or raw accuracy rather than operationally meaningful thresholds.
Our study addresses the latter two by fixing a daily alert quota and computing bootstrap
confidence intervals calibrated to the small count of available insider subjects.

\textbf{Li et al.}~\cite{li2021image} reframe log analysis as a vision problem: each
user's daily activity is encoded as a two-dimensional image and fed to a self-supervised
geometric-transformation classifier. Measured gains over a reconstruction baseline reach 4\,\%
at the session level and 2\,\% at the user level in AUROC.
AUROC aggregates performance over all operating thresholds and is therefore a poor proxy for
how many true insiders an analyst encounters when capped to a small daily queue.
The approach also offers no per-alert rationale.

\textbf{Tian et al.}~\cite{tian2025scenarios} observe that different insider archetypes leave
qualitatively different footprints and train dedicated multi-head-attention classifiers for each
scenario (privilege abuse, identity theft, data exfiltration), encoding behaviour as one-hot
daily sequences.
Relying on scenario-specific supervised models presupposes access to labelled attack traces
during training, a condition that rarely holds in live deployments.
Reported accuracy improvements of 2--5\,\% per scenario are difficult to interpret: at CERT's
positive rate of approximately 0.25\,\%, a trivial all-benign classifier already exceeds
99.75\,\% accuracy, so accuracy conveys essentially no information about detection utility.

\textbf{Kantchelian et al.}~\cite{kantchelian2024facade} introduce Facade, a contrastive
self-supervised model that evaluates each access request jointly against the requesting user's
role and the sensitivity of the target resource, trained on Google-internal document, SQL and
RPC access logs.
The reported operational false-positive rates are exceptionally low---under 0.01\,\% overall
and as low as three in a million for isolated rogue accesses---and this result directly
supports the design rationale behind our per-user baseline in C1: population-level anomaly
scores cannot achieve workable precision when user workloads are heterogeneous.
Two constraints limit comparability with our work: the evaluation data are proprietary and
cannot be independently reproduced, and alerts are not accompanied by human-readable
explanations, nor is any analyst study reported.

% ------------------------------------------------------------------
\subsection{Cloud Audit Log Anomaly Detection}
% ------------------------------------------------------------------

\textbf{Haunschild et al.}~\cite{haunschild2026devsecops} embed Isolation Forest inside a
DevSecOps monitoring pipeline for AWS CloudTrail logs and report $F_1 \approx 0.96$, claiming a
roughly 40\,\% relative gain over a prior threshold-based baseline.
This is the closest prior work to ours in the data source: it uses the same
cloud-audit-log structure and the same baseline detector (Isolation Forest) that we compare against.
Two gaps remain: the $F_1$ is reported at a single operating threshold, so it is unclear how
many insiders per day an analyst would find, and there is no explanation of why a user-day was
flagged.

\textbf{Zou et al.}~\cite{zou2025cloudanobench} introduce CloudAnoBench, a cloud anomaly
benchmark of 28 anomalous and 16 deceptive-normal scenarios (1,252 labelled cases, $\approx$200k
entries) and the LLM-agent detector CloudAnoAgent that couples symbolic verification with
language-model reasoning.
Their target is operational anomalies (faults, misconfigurations) rather than credential misuse
or insider threat; evaluation is per-case classification rather than a ranked daily queue.
Their characterisation of \emph{deceptive normal} behaviour---activity that looks legitimate but
is part of an attack---is directly relevant to the challenge we face in S3 (the disgruntled
administrator), and it motivates per-entity rather than population-level baselines.

% ------------------------------------------------------------------
\subsection{Conformal Prediction and Statistical Foundations}
% ------------------------------------------------------------------

\textbf{Bates et al.}~\cite{bates2023conformal} provide the theoretical grounding for using
split-conformal scores in outlier detection.
The central guarantee is a super-uniformity result: when a held-out calibration set is
exchangeable with a genuinely normal test observation, the empirical rank-based $p$-value
$(1 + \#\{c \ge s\}) / (n+1)$ is stochastically no larger than uniform under the null,
which means standard multiplicity corrections such as Benjamini--Hochberg FDR control remain
valid when many user-days are tested simultaneously.

\emph{Justification of the $p < 0.01$ secondary-family threshold.}
The threshold $p < 0.01$ that appears in the explanation line (``any other family with $p<0.01$'')
is not a detection threshold; it controls only which behaviour families are \emph{reported} alongside
the primary one.
We set it at 0.01 for three principled reasons, drawing directly on Bates et al.:

\begin{enumerate}
  \item \textbf{Calibration resolution.} Our pooled calibration set contains approximately $10^5$
    rows (all users' held-out days), so the minimum achievable $p$-value is
    $1/(n+1) \approx 10^{-5}$.
    A threshold of 0.01 is $1{,}000\times$ above the resolution floor, ensuring the threshold is
    never artefactually met because the calibration set is too small.
    By contrast, a per-user calibration set of 30 days has a resolution of $1/31 \approx 0.032$,
    which is already above 0.01; this is the quantitative reason pooling is necessary.
  \item \textbf{Conventional error control.} Bates et al.\ show that at threshold $\alpha$ the
    false-discovery rate of the outlier labels is controlled at $\alpha$ under exchangeability.
    We use $\alpha = 0.01$ (not the common 0.05) because the cost of a spurious secondary family
    appearing in an explanation is analyst confusion: the analyst would investigate a behaviour
    family that did not actually deviate.
    Setting $\alpha = 0.01$ means that, on a non-malicious day, the probability that any given
    non-primary family is incorrectly included in the explanation is at most 1\,\%.
    With four families and the primary already identified, the expected number of false secondary
    families per non-malicious alert is at most $3 \times 0.01 = 0.03$.
  \item \textbf{Empirical validation on the development period.}
    On the development-period calibration data we swept $\alpha \in \{0.001, 0.005, 0.01, 0.05\}$.
    At $\alpha = 0.05$ roughly 12\,\% of non-malicious alerts carried a spurious secondary family;
    at $\alpha = 0.01$ this dropped below 4\,\% while secondary families were still surfaced for
    the majority of true alerts where two behaviour families genuinely co-deviated
    (e.g.\ S1 shows both resource and temporal deviations on data-exfiltration days).
    $\alpha = 0.01$ was therefore fixed before the test window was scored.
\end{enumerate}

The detection ranking itself uses the combined $p$-value (C3, fusion step), not the 0.01 cutoff.
The threshold governs only explanation verbosity and is therefore a design decision about
analyst experience, grounded in the exchangeability guarantee of Bates et al.

% ------------------------------------------------------------------
\subsection{SOC Practice and Evaluation Methodology}
% ------------------------------------------------------------------

\textbf{AlAhmadi et al.}~\cite{alahmadi2022false} gather evidence from 20 survey respondents
and 21 semi-structured interviews with practising SOC analysts across several organisations.
A recurring finding is that the term ``false positive'' is applied by practitioners to events
that the detector correctly identified as anomalous but that turned out to be benign in
operational context---a labelling artefact rather than a detector failure.
Qualitative analysis distils five properties that distinguish actionable from unactionable alarms:
reliability, explainability, analytical depth, contextual grounding, and transferability across
teams (REACT).
This study motivates two specific design choices in our work: capping the daily queue at a size
an analyst can realistically review, and attaching a behaviour-family explanation to every alert.
AlAhmadi et al.\ neither build a detector nor report any budget-level performance numbers.

\textbf{Arp et al.}~\cite{arp2022dos} audit thirty papers published at top security venues
and catalogue ten recurring methodological weaknesses that systematically overstate detection
performance, among them temporal data leakage, reliance on aggregate AUC at very low base rates,
and failure to account for class imbalance when choosing metrics.
The analysis reveals that whole-curve summaries are particularly misleading when the fraction of
malicious instances is well below one percent, exactly the regime of our dataset.
We address this by fixing each evaluation at a concrete daily alert count rather than integrating
over thresholds, by ensuring all feature statistics use only past data, and by locking
hyperparameters before any test-period observations are processed.

% ------------------------------------------------------------------
\subsection{Explainability in Security ML and Alert Fatigue}
% ------------------------------------------------------------------

\textbf{Nadeem et al.}~\cite{nadeem2023sok} survey explainability practice across 118 security
ML papers. A striking finding is that fewer than one in seven include any form of user study,
and the dominant purpose of the explanation is to satisfy model developers rather than to assist
the practitioner who must act on each alert. Technical quality metrics such as attribution
fidelity and output stability are rarely connected to real analyst outcomes.
Objective O4 in our work is a direct response to this gap: instead of assuming that a
high-fidelity explanation is operationally useful, we directly measure the accuracy of four
human reviewers triaging alerts under an explanation-only condition against a score-only
condition in a pre-registered crossover design.

\textbf{Grzeczkowicz et al.}~\cite{grzeczkowicz2025xai} repurpose anomaly-detector SHAP
attributions as input features for a downstream classifier that predicts which insider \emph{scenario}
is active, distinguishing privilege abuse from data exfiltration on the CERT corpus.
The SHAP vectors are shown to retain scenario-discriminative structure even after aggregation
to the daily level, a finding that motivates our use of feature-level deviation magnitudes in
explanations. Neither a fixed-budget evaluation nor a human triage study is reported.

\textbf{Faqihi et al.}~\cite{faqihi2026shap} employ exact Shapley values not to support
human judgement but to automatically reassign detector output labels, targeting label correction
rather than analyst assistance.
Evaluation on 316{,}250 CERT instances reports accuracy and false-positive rates, but as
discussed above, accuracy at a sub-1\,\% positive rate carries little diagnostic value, and
no analyst is involved in the evaluation pipeline.

\textbf{Ndichu et al.}~\cite{ndichu2026pact} tackle alert fatigue at a different stage of the
pipeline: their PACT system places a frozen XGBoost classifier in front of the analyst queue
and triggers an active-learning update only when a distributional shift detector fires,
limiting annotation requests to genuinely ambiguous periods.
On two IDS alert corpora (AIT-ADS, BOTSv1) the method reduces the false-positive load by 43\,\%
and 21\,\% respectively while keeping recall within an acceptable window.
PACT receives already-scored IDS alerts as input and requires analyst-labelled feedback for
its retraining loop, which distinguishes it from our setting where alerts are generated from
raw audit logs and no labels are available during deployment.
The core quantitative insight---that suppressing false positives necessarily trades away some
true-positive coverage---is precisely the budget-precision trade-off that our DR@B and P@B
curves characterise across multiple daily alert quotas.

% ------------------------------------------------------------------
\subsection{Summary and Positioning}
% ------------------------------------------------------------------

Table~\ref{tab:litreview} summarises the fourteen works reviewed above against the four
evaluation dimensions that define our contribution.

\begin{table}[t]
\centering
\caption{Summary of related work (Y = yes, N = no, $\dagger$ = partial).
``Budget-level'' means the paper reports detection at a fixed daily or per-session analyst budget.}
\label{tab:litreview}
\resizebox{\columnwidth}{!}{%
\begin{tabular}{lp{3cm}cccc}
\toprule
Ref & Method & Budget eval & Expl. & Human study & Cloud logs \\
\midrule
\cite{le2021ensembles}         & Unsupervised ensemble  & Y$\dagger$ & N & N & N \\
\cite{yuan2021deep}            & Survey                 & N          & N & N & N \\
\cite{li2021image}             & Image + self-supervised& N          & N & N & N \\
\cite{tian2025scenarios}       & Multi-head attention   & N          & N & N & N \\
\cite{kantchelian2024facade}   & Contrastive per-user   & N          & N & N & N \\
\cite{haunschild2026devsecops} & Isolation Forest       & N          & N & N & Y \\
\cite{zou2025cloudanobench}    & LLM agent + benchmark  & N          & N & N & Y \\
\cite{bates2023conformal}      & Conformal $p$-values   & --         & --& --& -- \\
\cite{alahmadi2022false}       & SOC practitioner study & N          & --& Y & -- \\
\cite{arp2022dos}              & Pitfall survey         & N          & N & --& -- \\
\cite{nadeem2023sok}           & XAI survey             & N          & --& --& -- \\
\cite{grzeczkowicz2025xai}     & SHAP scenario classif. & N          & Y$\dagger$ & N & N \\
\cite{faqihi2026shap}          & SHAP re-classification & N          & Y$\dagger$ & N & N \\
\cite{ndichu2026pact}          & Active-learning triage & Y$\dagger$ & N & N & N \\
\midrule
\textbf{Ours}                  & BudgetAwareConformal   & \textbf{Y} & \textbf{Y} & \textbf{Y} & Y$\dagger$ \\
\bottomrule
\end{tabular}}
\end{table}

No prior work simultaneously evaluates at a fixed daily alert budget, provides a per-alert
explanation, and validates those explanations with human reviewers.
Le and Zincir-Heywood~\cite{le2021ensembles} and Ndichu et al.~\cite{ndichu2026pact} are the
closest on the budget dimension, but both use session-share or recall-within-window budgets rather
than a fixed analyst queue; neither produces explanations nor runs human studies.
Kantchelian et al.~\cite{kantchelian2024facade} achieve the strongest precision on proprietary
data but provide no explanation study.
Our cloud-log tick is marked partial because CERT r4.2 contains no cloud control-plane events;
we treat it as an audit log of the same structural type (who, what, where, when).
